% Reconstructed text from V5 PDF page 26; non-native source
\par\smallskip\noindent\rule{\textwidth}{0.4pt}\par{\sffamily\footnotesize Original V5 page 26 -- recovered text}\par\begin{flushleft}\footnotesize\ttfamily
\par\smallskip
\par\smallskip
\noindent Diagnostics and Serving Infrastructure for MAS 26\par
\par\smallskip
\par\smallskip
\par\smallskip
\noindent starvation under greedy scheduling (§8), restriction to a single memory tier (GPU only; ex-\par
\noindent tending to CPU/disk is future work), and exact-match-only prefix lookup (no fuzzy or semantic\par
\noindent matching).\par
\par\smallskip
\noindent Cost-benefit eviction. Mooncake [32] disaggregates KV cache from compute into a dis-\par
\noindent tributed cache pool called KVCache Pool, accessed via one-sided RDMA reads that bypass\par
\noindent the serving engine’s critical path. The architecture separates prefill nodes (which generate KV\par
\noindent blocks) from decode nodes (which consume them), with cache blocks stored on dedicated cache\par
\noindent nodes. Mooncake implements configurable eviction policies (LRU, LFU, or hybrid) with a cost-\par
\noindent benefit model that weighs cache retention value against RDMA transfer bandwidth. Under a\par
\noindent prefix-heavy workload, Mooncake achieves up to 525\% throughput improvement while main-\par
\noindent taining latency SLAs. The system received the FAST 2025 Best Paper Award. For multi-agent\par
\noindent scenarios, Mooncake’s disaggregated architecture decouples cache lifetime from compute life-\par
\noindent time: a cached prefix can persist on cache nodes even after the compute node that generated it\par
\noindent has moved to other requests. However, the eviction policies remain request-level (LRU/LFU)\par
\noindent and do not incorporate agent identity, workflow structure, or diagnostic value.\par
\noindent  InferCept [1] addresses a challenge specific to augmented LLM workloads: when generation\par
\noindent is intercepted for an external tool call, should the system preserve the in-progress KV cache\par
\noindent (wasting GPU memory while waiting) or discard it (wasting compute on later recomputation)?\par
\noindent InferCept introduces min-waste adaptive scheduling (§4.3, Equation 5) that makes a per-request\par
\noindent decision among three strategies: Preserve (retain in GPU memory), Swap (migrate to CPU\par
\noindent memory), or Discard (recompute upon resumption). The scheduling algorithm sorts intercepted\par
\noindent requests by expected GPU waste in descending order, assigns a per-iteration swap budget to\par
\noindent the highest-waste requests, and for the remainder selects whichever of Preserve or Discard\par
\noindent minimizes expected waste. Key quantitative results: on GPT-J-6B with a single A100, InferCept\par
\noindent sustains 1.6× higher request rate and reduces normalized latency by 1.9–5.7× (§5); on Llama3-\par
\noindent 70B across 4×A100, it achieves 2× higher throughput with 1.3–12× latency reduction. The\par
\noindent full optimization stack reduces GPU waste from 27\% (discard baseline) to just 0.69\% (§5.2,\par
\noindent Figure 3). InferCept’s duration estimation achieves 93\% of an oracle policy that knows the exact\par
\noindent tool-call duration in advance (§4.4). For cross-layer analysis, InferCept’s three-state model is\par
\noindent directly relevant: the Preserve/Swap/Discard decision is precisely the kind of retention trade-\par
\noindent off our evidence-loss formalization captures. The diagnostic question is whether the intercepted\par
\noindent state contains diagnostically valuable signals (tool-call context, intermediate reasoning) whose\par
\noindent loss during a Discard event constitutes an Li (T ) event (ICML 2024).\par
\noindent  Continuum [16] addresses the multi-turn agentic serving challenge: when an agent makes\par
\noindent a tool call, should the serving system retain its KV cache (occupying GPU memory during\par
\noindent the external call) or evict it (requiring recomputation when the agent resumes)? The system\par
\noindent answers this with a TTL-based cost-benefit optimization (§4.1): each request’s time-to-live τ ∗\par
\noindent is computed as the value that maximizes expected benefit, balancing the probability that the\par
\noindent tool call completes before the TTL expires (estimated from an empirical CDF of historical tool-\par
\noindent call durations) against the memory opportunity cost and prefill-reload cost (profiled offline per\par
\noindent hardware/model pair). The scheduling priority is a multi-key tuple: preempted requests first,\par
\noindent then pinned/within-TTL requests, then program-level FCFS by arrival time (§4.3). A cold-start\par
\noindent mechanism uses three tiers of CDF estimation based on the number of collected samples (§4.2).\par
\noindent On real SWE-agent workloads at a production deployment, Continuum achieves up to 8.18× job\par
\noindent completion time improvement (§6.4), with 50.31\% cost reduction and 27.28\% latency reduction\par
\noindent on average. On synthetic multi-turn benchmarks, throughput improvements range from 1.10×\par
\noindent to 3.22× and delay reductions from 1.12× to 3.66× across A100, B200, and H100 hardware\par
\end{flushleft}
